%! Characteristics of Linear Functions — Unit 1, Lesson 1.0 — Experience & Formalize
%! (Activity · QuickNotes · Application · Check Your Understanding) (Answer Key)
% Directory stays `experience/` (build identifier in shared/lesson.mk); the label the student
% and the teacher read is "Experience & Formalize".
% Math Medic sizing: 12pt body, OPEN white space after every question (no writing lines).
% \answerspace{H}{} reserves exactly H in the blank; the key passes the answer as the second
% argument, so the two files paginate identically by construction.
% Page budget (user, 2026-08-20): Activity 2pp max · QuickNotes 1/2pp · Application 1/2-1pp ·
% Check Your Understanding 1-2pp (practice, NOT scored — there is no homework).
\documentclass[12pt]{article}
\usepackage{algebra2-article}
\usepackage{algebra2-key}
\newcommand{\answerspace}[2]{\par\nopagebreak\noindent\begin{minipage}[t][#1][t]{\linewidth}%
  \color{keyred}\bfseries #2\end{minipage}\par}
\newcommand{\lgrid}[4]{%
  \draw[step=1, linegray!40, very thin] (#1,#3) grid (#2,#4);
  \draw[->, semithick, charcoal] (#1,0)--(#2,0) node[below right, font=\scriptsize]{$x$};
  \draw[->, semithick, charcoal] (0,#3)--(0,#4) node[left, font=\scriptsize]{$y$};}
% Temperature graph: hours h on the horizontal axis (0..7), temperature T on the vertical (-8..8).
% The h-axis is drawn at T=0 so the crossing is the point students circle.
\newcommand{\tempgrid}{%
  \draw[step=1, linegray!40, very thin] (0,-8) grid (7,8);
  \draw[->, semithick, charcoal] (0,0)--(7.6,0) node[right, font=\scriptsize]{$h$};
  \draw[->, semithick, charcoal] (0,-8.6)--(0,8.8) node[above, font=\scriptsize]{$T$ ($^\circ$C)};
  \foreach \x in {1,...,7} \node[charcoal, font=\tiny, below] at (\x,0){\x};
  \foreach \y in {-8,-6,-4,-2,2,4,6,8} \node[charcoal, font=\tiny, left] at (0,\y){\y};}

\begin{document}

\pageheader{Unit 1, Lesson 1.0}{Experience \& Formalize: Freezing Point --- Answer Key}

\vspace{0.06in}

\begin{headlinebox}{forestbg}
\textbf{Freezing Point.} On a cold weekend the number that matters most is \emph{zero}: above it the
roads are wet, below it they are ice. Work the two temperature stories below with your group using
what you already know --- tables, rules, points --- and be ready to explain \emph{how you know} each
answer \textbf{from the graph}.
\end{headlinebox}

\vspace{0.04in}

% ══ Part 1 of 4 — ACTIVITY (20 min, groups of 3–4, prior knowledge only) ══════
\begin{scenariobox}[1. Saturday Morning]{forest}
At 6:00 a.m. on Saturday the temperature in Blacksburg was $-6\,^\circ$C. It rose steadily,
$2\,^\circ$C every hour, all morning (until noon). Let $h$ be the number of hours after 6:00 a.m.
and $T$ the temperature in $^\circ$C.

\begin{enumerate}[label=\textbf{\alph*.}, leftmargin=1.8em, itemsep=5pt]
  \item Complete the table.
    \begin{center}
    \renewcommand{\arraystretch}{1.5}
    \begin{tabular}{|c|c|c|c|c|c|c|c|}
    \hline
    $h$ & $0$ & $1$ & $2$ & $3$ & $4$ & $5$ & $6$ \\ \hline
    $T$ & $-6$ & $-4$ & \ans{$-2$} & \ans{$0$} & \ans{$2$} & \ans{$4$} & \ans{$6$} \\ \hline
    \end{tabular}
    \end{center}
  \item Write a rule for the temperature $h$ hours after 6:00 a.m. \quad
        $T(h) = $ \ans{$2h - 6$}
\end{enumerate}

\begin{minipage}[t]{0.54\linewidth}
\vspace{0pt}
\begin{enumerate}[label=\textbf{\alph*.}, leftmargin=1.8em, itemsep=5pt]
  \item Write a rule: \; $T(h) = $ \ans{$6 - 2h$}
  \item The line crosses the $T$-axis at $({\color{keyred}\mathbf{0}},\ {\color{keyred}\mathbf{6}})$ and the $h$-axis
        at $({\color{keyred}\mathbf{3}},\ {\color{keyred}\mathbf{0}})$. What does each crossing tell you about Sunday
        evening?
        \answerspace{2.4cm}{$(0,6)$: at 4:00 p.m.\ it was $6\,^\circ$C --- the starting temperature. $(3,0)$: three hours later, at 7:00 p.m., it reached $0\,^\circ$C --- freezing.}
\end{enumerate}
\end{minipage}\hfill
\begin{minipage}[t]{0.42\linewidth}
\vspace{0pt}\centering
\begin{tikzpicture}[x=0.58cm, y=0.26cm]
  \tempgrid
  \draw[very thick, forest] (0,6)--(7,-8);
  \fill[charcoal] (0,6) circle (2.8pt);
  \fill[charcoal] (3,0) circle (2.8pt);
\end{tikzpicture}
\end{minipage}

\begin{enumerate}[label=\textbf{\alph*.}, start=3, leftmargin=1.8em, itemsep=5pt]
  \item At 5:00 p.m. ($h=1$) the temperature is \ans{$4$} $^\circ$C. Is that number
        \textbf{positive} or \textbf{negative}? At that moment, is the temperature \textbf{rising}
        or \textbf{falling}? \\[4pt]
        So: can a temperature be \emph{positive} and \emph{falling} at the same time? Explain
        what each word is describing.
        \answerspace{2.2cm}{$4$ is positive, and the temperature is falling. Yes --- both at once. ``Positive'' describes the \emph{value} (above zero, graph above the $h$-axis); ``falling'' describes the \emph{direction} the values are moving (the line goes down left to right).}
  \item Above freezing: $T>0$ when \ans{$0\le h<3$} \qquad
        Below freezing: $T<0$ when \ans{$h>3$}
  \item \textbf{The story, then the math.} In the story, which hours actually make sense as
        inputs? \ans{$h\ge0$, the evening hours} \\[2pt]
        Now \emph{forget the story}: the rule $T(h) = 6 - 2h$ accepts \emph{any} number $h$ and its
        line runs forever both ways, so all the possible inputs are \ans{all reals} and all
        the possible outputs are \ans{all reals}. Why are those two answers different?
        \answerspace{1.5cm}{The rule allows every real number, but the \emph{story} only starts at 4:00 p.m., so negative $h$ describes nothing --- a story keeps only the inputs that make sense.}
\end{enumerate}
\end{scenariobox}

% ══ Part 2 of 4 — QUICKNOTES (13 min debrief fills this; target half a page) ══
\boxguard[14]
\begin{tcolorbox}[enhanced, colback=sky, colframe=navy, boxrule=0.8pt, arc=2mm,
    left=3mm, right=3mm, top=1.5mm, bottom=1.5mm, breakable,
    fonttitle=\bfseries\small\color{white}, title={QuickNotes: Characteristics of a Linear Function},
    attach boxed title to top left={yshift=-2mm, xshift=4mm},
    boxed title style={colback=navy, colframe=navy, arc=1.5mm}]
\small
\begin{minipage}[t]{0.31\linewidth}
\vspace{2pt}\centering
\begin{tikzpicture}[scale=0.26]
  \lgrid{-2}{6}{-6}{6}
  \draw[very thick, forest] (-1,-6)--(5,6);
  \draw[semithick, navy, dashed] (2,0)--(3,0)--(3,2);
  \node[navy, font=\tiny] at (2.5,-0.7){run $1$};
  \node[navy, font=\tiny, right] at (3.05,1){rise $2$};
  \fill[charcoal] (0,-4) circle (3.4pt);
  \fill[charcoal] (2,0) circle (3.4pt);
\end{tikzpicture}\\[1pt]
{\footnotesize Example: $f(x) = 2x - 4$}
\end{minipage}\hfill
\begin{minipage}[t]{0.65\linewidth}
\vspace{0pt}
\begin{itemize}[leftmargin=1.1em, itemsep=3pt, topsep=0pt]
  \item A \textbf{linear function} has a \ans{constant} rate of change, so its graph is a
        \ans{straight line}. Form: $f(x) = mx + b$.
  \item \textbf{Slope} $m = \dfrac{\Delta y}{\Delta x} = \dfrac{\text{rise}}{\text{run}}$. \;
        For $f$: $m = {\color{keyred}\mathbf{2}}$.
  \item \textbf{$y$-intercept} $(0,b)$: set $x=0$. For $f$: $(0,{\color{keyred}\mathbf{-4}})$. \quad
        \textbf{$x$-intercept / zero} $(a,0)$: solve $f(x)=0$. For $f$: $({\color{keyred}\mathbf{2}},0)$.
\end{itemize}
\end{minipage}

\vspace{4pt}
\begin{itemize}[leftmargin=1.1em, itemsep=3pt, topsep=0pt]
  \item \textbf{Increasing / decreasing} --- read left to right: \;
        \mbox{$m>0$: \ans{increasing}} \quad \mbox{$m<0$: \ans{decreasing}} \quad
        \mbox{$m=0$: \ans{constant}}
  \item \textbf{Positive / negative:} $f(x)>0$ where the graph is \ans{above} the $x$-axis;
        $f(x)<0$ where it is \ans{below}. The sign can only switch at a \ans{zero}.
        For $f$: positive when $x {\color{keyred}\mathbf{>2}}$, negative when $x {\color{keyred}\mathbf{<2}}$.
        \emph{Where it sits $\ne$ which way it moves} --- Sunday's line was positive \emph{and} falling.
  \item \textbf{Domain / range} of a line that is not horizontal: \ans{all reals} /
        \ans{all reals}. A horizontal line $f(x)=b$ has range \ans{$\{b\}$} and no
        zero (unless $b=0$). In a \emph{story}, the domain is only the inputs that make sense.
\end{itemize}
\end{tcolorbox}

% ══ Part 3 of 4 — APPLICATION (7 min, worked together; the modeling standard) ══
\boxguard[14]
\begin{notesbox}{Application: The Car-Wash Card}
We will do this one together. A car-wash gift card starts with \$40, and each wash costs \$8, so
the balance after $w$ washes is $B(w) = 40 - 8w$.

\begin{enumerate}[leftmargin=1.7em, itemsep=5pt]
  \item What does the slope $-8$ tell you about the card? And the $y$-intercept $(0,40)$?
    \answerspace{1.4cm}{Slope $-8$: every wash takes \$8 off the balance (the balance drops \$8 per wash). $(0,40)$: before any washes, the card holds \$40 --- the starting balance.}
  \item Find the zero of $B$ by solving $B(w) = 0$.
    \begin{work}
      40 - 8w &= 0 \\
          -8w &= -40 \\
            w &= 5
    \end{work}
    In the story, that zero means:
    \answerspace{1.4cm}{After 5 washes the card is empty --- the balance hits \$0 exactly at the 5th wash.}
  \item Which inputs $w$ make sense? \ans{whole $w$, $0$ to $5$} \;
        Is $B$ increasing or decreasing there? \ans{decreasing}
  \item Suppose each wash cost \$10 instead of \$8. Would the card run out \emph{sooner} or
        \emph{later}? \ans{sooner ($w=4$)} \; And on the graph, what changes --- where the
        line \emph{starts}, how \emph{steep} it is, or both? Say why.
    \answerspace{1.4cm}{Only the steepness. The card still starts at \$40, so the $y$-intercept is unchanged; the slope goes from $-8$ to $-10$, a steeper drop, so the line reaches zero sooner.}
\end{enumerate}
\end{notesbox}

% ══ Part 4 of 4 — CHECK YOUR UNDERSTANDING (10 min, practice, NOT scored) ═════
% This is the lesson's only practice set: there is no homework, so the six items
% below carry the whole A2.F.2 spread — graph, contrast pair, table, the constant
% special case with its boundary, a model in a fresh context, and the formative MC.
\boxguard[14]
\begin{notesbox}{Check Your Understanding \quad {\normalfont\itshape (practice --- not scored)}}
Six exercises --- the lesson's practice, and there is \textbf{no homework}. Work 1--3 with a partner,
then 4--6 on your own. For every answer, be ready to say \emph{how you know} it.

\begin{enumerate}[leftmargin=1.7em, itemsep=4pt]
  \item \textbf{Read the rule.} Complete the feature list for $f(x) = 2x + 2$.
    \begin{center}
    \renewcommand{\arraystretch}{1.4}
    \begin{tabular}{@{}l l l@{}}
      Slope: \ans{$2$} & $y$-intercept: \ans{$(0,2)$} & Zero: \ans{$x=-1$} \\
      Domain: \ans{all reals} & Range: \ans{all reals} & Incr.\ or decr.? \ans{increasing} \\
      Positive when $x$ \ans{$>-1$} & Negative when $x$ \ans{$<-1$} & Sign flips at: \ans{the zero}
    \end{tabular}
    \end{center}

  \boxguard[12]
  \item \textbf{Same zero, opposite behavior.} Both lines have their zero at $x = 1$.
    \begin{center}
    \begin{minipage}[c]{0.34\linewidth}\centering
    \begin{tikzpicture}[scale=0.36]
      \lgrid{-3}{5}{-4}{4}
      \draw[very thick, forest] (-2.5,-3.5)--(4.5,3.5) node[above left, font=\scriptsize, forest]{$f$};
      \draw[very thick, navy] (-2.5,3.5)--(4.5,-3.5) node[below left, font=\scriptsize, navy]{$g$};
      \fill[charcoal] (1,0) circle (3.4pt);
    \end{tikzpicture}
    \end{minipage}\hfill
    \begin{minipage}[c]{0.62\linewidth}
    \begin{enumerate}[label=(\alph*), leftmargin=1.9em, itemsep=4pt]
      \item Increasing: \ans{$f$} \; Decreasing: \ans{$g$}
      \item $f(x)>0$ when $x$ \ans{$>1$} \quad $g(x)>0$ when $x$ \ans{$<1$}
      \item At $x=3$: $f$ is \ans{positive} \; $g$ is \ans{negative}
    \end{enumerate}
    \end{minipage}
    \end{center}
    (d) Why can a line's sign change at its \emph{zero} and nowhere else?
    \answerspace{1.4cm}{The zero is the only place the line touches the $x$-axis. To get from below the axis to above it, the graph has to cross --- and a non-horizontal line crosses exactly once.}

  \item \textbf{From a table.} A linear function is given below. \emph{Watch the step in $x$.}
    \begin{center}
    \renewcommand{\arraystretch}{1.3}
    \begin{tabular}{|c|c|c|c|c|}
    \hline
    $x$ & $0$ & $2$ & $4$ & $6$ \\ \hline
    $f(x)$ & $-6$ & $-3$ & $0$ & $3$ \\ \hline
    \end{tabular}
    \end{center}
    Slope $m = \dfrac{\Delta y}{\Delta x} = {\color{keyred}\mathbf{\tfrac{3}{2}}}$ \quad
    $y$-intercept: \ans{$(0,-6)$} \quad Zero: \ans{$x=4$} \\[4pt]
    Increasing or decreasing? \ans{increasing} \quad $f(x)<0$ when $x$ \ans{$<4$}

  \boxguard[12]
  \item \textbf{The flat one.} The graph of $c(x) = -2$ is a horizontal line through $(0,-2)$.
    \begin{enumerate}[label=(\alph*), leftmargin=1.9em, itemsep=4pt]
      \item Slope: \ans{$0$} \; Incr., decr., or constant? \ans{constant} \;
            Range: \ans{$\{-2\}$} \; A zero? \ans{no}
      \item Is $c(x)$ positive or negative? \ans{negative everywhere}
      \item For which value of $b$ does $c(x) = b$ \emph{have} a zero? \ans{$b=0$} \;
            And then how many zeros does it have? \ans{every input is one}
    \end{enumerate}

  \boxguard[12]
  \item \textbf{Model it.} A pool already holds $600$ gallons when a hose is turned on; the hose
        adds $150$ gallons a minute, so $W(t) = 600 + 150t$ after $t$ minutes.
    \begin{enumerate}[label=(\alph*), leftmargin=1.9em, itemsep=4pt]
      \item Slope $150$ means \ans{$+150$ gal each minute} \;
            $(0,600)$ means \ans{600 gal at $t=0$}
      \item Find the zero of $W$ by solving $W(t) = 0$.
        \begin{work}
          600 + 150t &= 0 \\
                150t &= -600 \\
                   t &= -4
        \end{work}
        Does that zero mean anything in this story? Explain.
        \answerspace{1.2cm}{No. $t=-4$ is four minutes \emph{before} the hose was turned on, and the story only starts at $t=0$ --- the model does not describe negative time.}
      \item Which inputs $t$ make sense? \ans{$t\ge 0$} \; Is $W$ increasing or
            decreasing on them? \ans{increasing}
    \end{enumerate}

  \item \textbf{Multiple choice.} For the linear function $y = -3x + 6$, which statement is
        \emph{true}?
    \begin{enumerate}[label=\Alph*., leftmargin=2.2em, itemsep=1pt]
      \item[\textbf{A.}] The function is increasing. \hfill
            \textbf{B.} It is positive for all $x > 2$.
      \item[\textbf{C.}] The zero is at $x = 2$. \textcolor{keyred}{\textbf{$\leftarrow$ correct}} \hfill
            \textbf{D.} The $y$-intercept is $(0,-3)$.
    \end{enumerate}
    Name one of the other three and say why it is false.
    \answerspace{1.4cm}{C is true: $-3x+6=0$ gives $x=2$. (A: the slope $-3$ is negative, so it decreases. B: it is positive for $x<2$, negative for $x>2$. D: the $y$-intercept is $(0,6)$; $-3$ is the slope.)}
\end{enumerate}
\end{notesbox}

\boxguard[6]
\begin{spiralbox}
\textbf{Coming next --- Lesson 1.1: Linear Functions.} Today you \emph{read} a line's characteristics
off its graph, table, and rule. Next you run it backwards: slope and a point \emph{write} the
equation three ways, and shifting the parent line $y=x$ graphs it.
\end{spiralbox}

\end{document}
